\hypertarget{o-pauxeds-do-agro-uxe9-o-pauxeds-da-fome}{%
\section{202205281121 ─ O país do Agro é o país da
fome}\label{o-pauxeds-do-agro-uxe9-o-pauxeds-da-fome}}

\begin{quote}
``a realidade é que o \textit{agribusiness}, longe de ser a solução,
apenas agrava o problema da fome. Isso porque dele advém não apenas a
modernização da agricultura, mas a transferência de um modelo de
desenvolvimento econômico e de relações sociais para o Terceiro Mundo ─
o modelo capitalista de produção. Dessa forma, o \textit{agribusiness}
apenas exacerba as desigualdades sociais, que (\ldots) são as causas
reais da fome'' (\textit{Agribusiness in the Americas}, Burbach \&
Flynn, apud \textbf{Formação Política do Agronegócio}, Caio Pompeia,
p.~79)
\end{quote}

\hypertarget{introduuxe7uxe3o}{%
\subsection{Introdução}\label{introduuxe7uxe3o}}

Fazer uma discussão sobre o papel de agrotóxicos no tocante a qualquer
faceta social no Brasil ─ em particular à fome ─ faz necessário uma
discussão sobre o avanço do agronegócio, e em particular sobre sua
influência disruptora sobre os sistemas sociais relacionados ─ direta ou
indiretamente.

Faz-se necessário, posto que o problema não é tão simples e imediato
quanto ``agrotóxicos causam fome''; temos graus diferentes tanto da
influência dos agentes envolvidos quanto das consequências engendradas,
posto que a discussão não envolve somente o uso de agrotóxicos, mas
também a ubiquidade de seu uso, o aparato legal que o permite, as
influências econômicas (e políticas) que as empresas que produzem tais
agrotóxicos exercem, etc. ─ e isso focando somente na questão
``agrotóxicos''; o buraco é bem mais embaixo. A discussão é complexa, em
linha com a visão marxista da totalidade concreta. E, como marxistas,
devemos nos colocar neste espaço mental a fim de criar uma discussão
efetiva sobre tal problema social.

\hypertarget{motivauxe7uxe3o}{%
\subsection{Motivação}\label{motivauxe7uxe3o}}

Atualmente há um ataque em várias frentes contra a população brasileira:
o que alguns chamaram de Pacote da Destruição. São projetos de lei (PLs)
impulsionados pela Bancada Ruralista\footnote{Formalmente chamada de FPA
  (Frente Parlamentar da Agropecuária).}, e que buscam legalizar atos
que são inconstitucionais, mas que, independente ou não da legalidade,
são contra os interesses da população e do meio ambiente.

!Pasted image 20220528141821.png (\textit{Fonte}:
\href{https://deolhonosruralistas.com.br/2022/02/10/passando-a-boiada-12-das-45-prioridades-do-governo-no-congresso-sao-no-campo/?utm_source=rss\&utm_medium=rss\&utm_campaign=passando-a-boiada-12-das-45-prioridades-do-governo-no-congresso-sao-no-campo}{Passando
a Boiada: 12 das 45 prioridades do governo no Congresso são no campo
(\textbf{De Olho nos Ruralistas})})

!Pasted image 20220605222817.png (\textit{Fonte:}
\href{https://deolhonosruralistas.com.br/2022/02/07/em-meio-a-avalanche-ruralista-maioria-dos-partidos-ignora-temas-agrarios/?utm_source=pocket_mylist}{Em
meio a avalanche ruralista, maioria dos partidos ignora temas agrários
(\textbf{De Olho nos Ruralistas})})

Os principais PLs desse Pacote são: - \textbf{PL 6.299/2002}:
``\textbf{PL do Veneno}''. Aprovado pela Câmara dos Deputados em
fevereiro de 2022 (301 a favor, 150 contra); foi encaminhado ao Senado,
onde tramita como \textbf{PL 526/1999}. Principais pontos: 1. O controle
da autorização de novos agrotóxicos deixa de ser feito pelo Ministério
da Saúde, Anvisa\footnote{Agência Nacional de Vigilância Sanitária},
Ministério do Meio Ambiente e Ibama, e começa a ser comandado pelo
\textbf{MAPA}\footnote{Ministério da Agricultura, Pecuária e
  Abastecimento}. * Os órgãos anteriormente responsáveis podem até
emitir pareceres e alertas de risco, mas o parecer final é dado
diretamente pelo MAPA 2. Alterar o termo ``agrotóxicos'' por *
\textbf{“pesticidas”} (``conforme'' a tendência mundial de chamá-los de
\textit{pesticides}) quando aplicados em plantações *
\textbf{“produtos de controle ambiental”} quando aplicados em florestas
e ambientes hídricos 3. Estimula prazos mais rápidos para registros de
agrotóxicos por órgãos federais (até dois anos) * Mesmo quando não
houver manifestação conclusiva dentro do prazo estipulado de sua
verificação, ele recebe uma \textbf{autorização temporária} para atuar
em território brasileiro * \textbf{PL 2.633/2020}:
``\textbf{PL da Grilagem}''. Aprovado pela Câmara e encaminhado ao
Senado em agosto de 2021 (aguarda votação da Comissão do Meio Ambiente).
* \textbf{PL 191/2020}: \textbf{Mineração em Terras Indígenas}. Ganhou
ímpeto após a deflagração da guerra na Ucrânia, sob a retórica de
produção nacional de fertilizantes frente ao seu aumento dos preços
internacionais (Rússia está entre nossos principais fornecedores de
fertilizantes) * \textbf{PL 490/2007}: ``\textbf{Marco Temporal}''.
Restringe e revoga severamente a demarcação de terras indígenas. \#\# Os
níveis de abstração da discussão \#\#\# Oposição entre capital e força
de trabalho No nível mais \textit{abstrato}\footnote{O que quer dizer:
  no nível mais ``simplificado'', desconsiderando fatores particulares e
  específicos.}, temos que o problema envolve a clássica oposição entre
capital e trabalho, posto que estamos em uma sociedade cujo modo de
produção se assenta sobre o trabalho assalariado\footnote{Ou seja,
  trabalhadores vendem sua força de trabalho, posto que não têm nenhum
  outro modo de ganhar seu sustento, e vendem-na para detentores de
  capital, os quais detêm a propriedade de meios de produção de maneira
  \textit{privada} (i.e.~excludente de outrem), e obtêm tanto seus
  lucros, quanto seu ``capital'' (no sentido cotidiano/vulgar) para
  reinvestimento em produções futuras, a partir do trabalho
  não-remunerado (mais-valia, ou mais-valor) dos trabalhadores que
  contratou em sua empresa.}, e em escala mundial. Porém, isso ainda não
nos diz nada sobre o problema específico que temos.

Porém, já deste nível conseguimos ter uma visão cristalina sobre o
agronegócio: como seu nome diz, ele é um \textit{negócio}, ele busca
produzir \textit{mercadorias} (\textit{commodities} em inglês). Ele
produz \textit{para vender}, e não \textit{para alimentar} (como somos
induzidos a pensar pelos meios ideológicos hegemônicos).

\begin{quote}
``Aqui {[}no processo capitalista de produção em geral{]}, os
\textit{valores de uso} só são produzidos
\textbf{porque e na medida em que} são o \textbf{substrato material}, os
\textbf{suportes} do \textit{valor de troca}. E, para nosso capitalista,
trata-se de duas coisas. Primeiramente, ele quer produzir um valor de
uso \textbf{que tenha um valor de troca} (\ldots). Em segundo lugar,
quer produzir uma mercadoria cujo valor seja
\textbf{maior do que a soma do valor das mercadorias requeridas para sua produção},
os meios de produção e a força de trabalho, para cuja compra ele
adiantou seu dinheiro no mercado.

Ele quer produzir não só um valor de uso, mas uma \textbf{mercadoria};
não só valor de uso, mas valor, e não só valor, mas também
\textbf{mais-valor}.'' (O Capital I, p.~263; grifo meu)
\end{quote}

Neste processo, podemos ver vários níveis\footnote{Podemos ver, por
  exemplo, a oposição entre os trabalhadores que trabalham nas empresas
  que produzem agrotóxicos, opostos ao capital da indústria química,
  embora esse nível não seja particularmente pertinente à nossa
  discussão. Faz-se necessário escolher algum ``nível'' a partir do qual
  começar a discussão, cuja ``distância'' do problema em questão ajude a
  ilustrá-lo melhor, invés de nos deixar com a pergunta ``tá, mas por
  que começar \textit{daí} e não \textit{de lá}?''.} dessa oposição.

As influências disruptivas do agronegócio sobre a infraestrutura social
de cidades e regiões inteiras, centralizando a atividade econômica
inteira destas em torno de \textit{suas} atividades, empregam a força de
trabalho local em - plantio de \textit{commodities} (soja, milho, trigo,
etc.), seja envolvendo trabalho ``braçal'' (manejando
tratores/colheitadeiras, atuando como jagunços de latifundiários\ldots{}
etc.) ou ``espiritual''/``mental'' (agrônomos, engenheiros,
meteorologistas, analistas de dados etc.\footnote{E a produção cultural
  que lhe dê sustentação, a versão midiática hegemônica de ``agro é
  pop'', etc.}) - transporte (circulação) das \textit{commodities}
produzidas (principalmente caminhoneiros, mas também condutores de
trem/barco etc.) - mineração, dada tanto a dependência do agronegócio em
fertilizantes (nitrogênio, potássio, fósforo), quanto os efeitos da
centralização na atividade mineradora na mitigação, em particular, da
agricultura familiar (e avanço consequente do agronegócio em terras que
ficaram ``avulsas'') - \textgreater{} ``Os capitalistas burgueses
favoreceram a operação
{[}\textit{de cercamentos na Inglaterra do século XVII, expulsão em massa de camponeses para liberação de terras para pastagens}{]},
entre outros motivos, para transformar o solo em artigo puramente
comercial, ampliar a superfície da grande exploração agrícola, aumentar
a oferta de proletários absolutamente livres, provenientes do campo
{[}\textit{i.e. expulsos de suas terras, forçados a ter de vender sua força de trabalho para sobreviver}{]}
etc.'' (O Capital, p.~796; destaques meus) - \textgreater{} ``Logo que o
trabalho começa a ser distribuído
{[}\textit{”forçosamente”, sem alternativas}{]}, cada um passa a ter um
campo de atividade exclusivo e determinado, que lhe é imposto e ao qual
não pode escapar; o indivíduo é caçador, pescador, pastor ou crítico
crítico, e assim deve permanecer se não quiser perder seu meio de vida''
(A Ideologia Alemã, p.~37-38; destaque meu) - Note-se que, aqui, a
influência do capital minerador com relação ao agronegócio não é
imediata, posto que a grande maioria dos fertilizantes empregados são
importados; ambos estes negócios buscam a exportação de seus produtos, e
somente pela importação do produto final (i.e.~o fertilizante) pode-se
perceber a verdadeira relação entre ambos (em particular, isso
desmistifica a retórica da PL da Mineração, de ``incentivar o mercado
interno de fertilizantes'' face ao pico nos preços de importação)

Em todo caso, há uma centralização da atividade econômica (ou seja, dos
trabalhos assalariados) da cidade em torno do agro(minero)negócio; este
finca suas raízes nas cidades em que age e suga todos os nutrientes que
pode, todos os meios de produção e força de trabalho a seu alcance,
empregando homens, leis e armas para assegurar seu próprio sustento; por
consequência, menos pessoas trabalhando em agricultura familiar (posto
que não paga mais as contas), que implica em menos alimentos produzidos
para a população (e a preços ``pagáveis'', destaque-se), e, por fim, um
alastramento dos empregos precarizados, da fome e da pobreza. Tal
centralização não é nada novo; é o clássico ``se não pode contra eles,
una-se a eles'', e como Marx disse: \textgreater{} ``O
{[}capitalista{]}, com um ar de importância, confiante e ávido por
negócios; o {[}assalariado{]}, tímido e hesitante, como alguém que
trouxe sua própria pele ao mercado e, agora, não tem nada mais a esperar
além da\ldots{} esfola.'' (O Capital I, p.~251).

\hypertarget{a-fome-uxe9-um-fenuxf4meno-social-poluxedtico}{%
\paragraph{A fome é um fenômeno
social-político}\label{a-fome-uxe9-um-fenuxf4meno-social-poluxedtico}}

Visto neste nível de abstração, pode-se ter uma visão mais abrangente
também sobre o surgimento do fenômeno da fome dentro do capitalismo: não
obstante ter produzido um ``mundo de riqueza e de cultura'', com uma
força produtiva mais do que capaz de produzir o suficiente para
todos\footnote{No tocante tanto às suas ``necessidades bíblicas'', como
  Löwy o põe em
  ``\href{https://www.ifch.unicamp.br/criticamarxista/arquivos_biblioteca/artigo164artigo2.pdf}{\textit{Ecossocialismo e planejamento democrático}}''
  (``água, comida, roupa, habitação''), quanto também suas necessidades
  social-culturais (lazer, entretenimento, socialização, comensalidade,
  etc.).}, vemos uma exacerbação da acumulação do capital nas mãos de
poucos, seja no baixo número de megacorporações e de bilionários (os
quais sempre sofrem alterações mediante crises que causam falências e
\textit{merges} de companhias), seja na centralização da atividade
econômica de cidades em torno dos interesses de empresas privadas (em
nosso caso, quase sempre multinacionais). Este último ponto merece um
destaque maior, dado sua ubiquidade no território brasileiro.

\st{Imaginemos alguma empresa multinacional chegando em território
brasileiro, em uma bela cidadezinha no meio do nada, longe de centros
urbanos e dos olhos da mídia estadual/nacional. Chegando ali, ela deseja
contratar mão-de-obra o mais barato possível, cuja melhor aposta seria
se absolutamente \textit{todos os moradores} desejassem trabalhar para
ela (a máxima oferta de trabalho possível, consequentemente pelo menor
preço salarial possível\footnote{Pelo conceito de ``exército industrial
  de reserva''.}). Porém, não só de comércio cria-se capital; também
através de acordos (e subornos) políticos consegue o capital assegurar
sua dominação sobre uma cidade.}

\hypertarget{oposiuxe7uxe3o-entre-monocultura-e-agroecologia}{%
\subsubsection{Oposição entre monocultura e
agroecologia}\label{oposiuxe7uxe3o-entre-monocultura-e-agroecologia}}

Da mesma forma que ninguém se atreveria a morar em um edifício com um só
alicerce, os ecossistemas mais estáveis na natureza são aqueles que
possuem a maior diversidade de espécies que o compõem.

Isso é pois tais ecossistemas são o que cientistas chamam de ``sistemas
complexos'', onde ``o todo é maior que a soma das partes''; vendo de um
ponto de vista dialético, percebe-se que, mesmo havendo uma clara
oposição entre as espécies componentes do sistema (o consumo dos
recursos locais por uma delas implica no não-consumo ─ ao menos
imediatamente ─ desses recursos por outras), a prevalência de uma sobre
todas as outras causa uma disrupção no sistema, o que retroage na
diminuição da população desta espécie, trazendo o sistema de volta à
homeostase. Ou seja, mesmo as espécies efetivamente buscando sua própria
sobrevivência, elas acabam dependendo umas das outras, perecendo caso
não haja o movimento recíproco de outras partes.

Tais ecossistemas são mais robustos que, por exemplo, suas ``partes
componentes'' vistas isoladamente, pois ocorre uma ``economia circular''
dentro deles: os resíduos de uma espécie servem como nutrientes para
outras espécies, e assim vai, num ciclo virtuoso; as necessidades de
certas partes da totalidade são satisfeitas dentro dela própria, por
outras componentes. Além disso, claro, o processo de seleção natural que
ocorreu (e ocorre) há centenas de milhões de anos, restando hoje os
sistemas melhor adaptados a seus ambientes e que possuam a melhor
sinergia interna.

A agroecologia busca ser um movimento de retorno à harmonia dos sistemas
naturais; mesmo sendo uma atividade que busque ser uma fonte de recursos
para o homem, ela busca colocar o homem \textbf{dentro} do ecossistema
(ou melhor, relembrá-lo disso), invés de somente usurpá-lo sem dar algo
de volta, como se fosse algo \textbf{externo}.

!Pasted image 20220530191942.png (Fonte:
\href{https://tcpermaculture.com/site/2013/05/27/nine-layers-of-the-edible-forest-garden/\#prettyPhoto}{Temperate
Climate Permaculture})

Sob essa perspectiva, torna-se óbvio a destruição que monoculturas
trazem (e que trazem ``para si próprias'', em boa medida): - a
necessidade constante de fertilizantes, posto que ``não é permitido''
que outras espécies ─ produtoras de nutrientes que poderiam ser-lhe
úteis, e às quais poderia também auxiliar ─ convivam no mesmo lote de
terra - a intensidade das secas, posto que\footnote{Dentre vários
  motivos, como o próprio aquecimento global, temperaturas maiores e a
  disrupção consequente de ecossistemas, etc.} não existem plantas que
sirvam de quebra-vento; portanto, os ventos fortes acabam levando
consigo a umidade do ar e dos estômatos das folhas - a hecatombe que
pestes têm o potencial de causar sobre safras inteiras, posto que não
existem insetos ou microorganismos que possam combater essas doenças

Desse ponto de vista, pode-se dizer que as monoculturas têm toda a terra
de sua propriedade ao seu alcance, mas a que custo? Prevaleceram somente
elas, com a mão férrea do homem a seu dispor, mas na natureza
prevaleceram justamente aquelas espécies que aprenderam a sobreviver
juntas entre si. É somente do ponto de vista financeiro que faz sentido
a ideia de que ``produzir mais é melhor''; na natureza isso é suicida. E
inevitavelmente traz a decadência àqueles que têm sede demais ao pote.

\st{Da mesma forma que o capitalista consegue ganhos ``extraordinários''
quando primeiro tem acesso à máquina em seu processo de produção, mas
logo percebe que o mais-valor vem somente da parcela variável de seu
capital (i.e.~da força de trabalho empregada); mutatis mutandis quanto à
alta produtividade de uma monocultura em sua primeira safra, e seu
decaimento em performance ao longo do tempo, precisando ser compensado ─
surprise, surprise ─ pela intervenção humana.}

\hypertarget{oposiuxe7uxe3o-entre-latifuxfandio-e-agricultura-familiar}{%
\subsubsection{Oposição entre latifúndio e agricultura
familiar}\label{oposiuxe7uxe3o-entre-latifuxfandio-e-agricultura-familiar}}

\hypertarget{sobre-a-questuxe3o-da-fome}{%
\subsection{Sobre a questão da fome}\label{sobre-a-questuxe3o-da-fome}}

\begin{quote}
``\ldots{} devemos começar por constatar o primeiro pressuposto de toda
a existência humana e também, portanto, de toda a história, a saber, o
pressuposto de que os homens têm de \textbf{estar em condições de viver}
para poder `fazer história'. Mas, para viver, precisa-se, antes de tudo,
de comida, bebida, moradia, vestimenta e algumas coisas mais. O primeiro
ato histórico é, pois,
\textbf{a produção dos meios para a satisfação dessas necessidades}, a
produção da própria vida material, e este é, sem dúvida, um ato
histórico, uma condição fundamental de toda a história, que ainda hoje,
assim como há milênios, tem de ser cumprida diariamente, a cada hora,
simplesmente para manter os homens vivos.'' (A Ideologia Alemã, p.~33;
grifo meu)
\end{quote}

\hypertarget{soberania-e-seguranuxe7a-alimentar-e-nutricional-ssan}{%
\subsubsection{Soberania e segurança alimentar e nutricional
(SSAN)}\label{soberania-e-seguranuxe7a-alimentar-e-nutricional-ssan}}

\begin{quote}
``1. A soberania alimentar é o caminho para erradicar a fome e a
subnutrição, e garantir a segurança alimentar duradoura e sustentável
para todos os povos. Entendemos por soberania alimentar o direito dos
povos a definir suas \textbf{próprias} políticas e estratégias
\textbf{sustentáveis} de produção, distribuição e consumo de alimentos,
que garantam o direito à alimentação a toda a população, com base na
\textbf{pequena e média produção}, respeitando suas
\textbf{próprias culturas} e diversidade dos modos camponeses de
produção, de comercialização e de gestão, nos quais a mulher desempenha
um papel fundamental'' (Fórum Mundial de Soberania Alimentar ─ Havana,
2001; grifo meu)

``a realidade é que o \textit{agribusiness}, longe de ser a solução,
apenas agrava o problema da fome. Isso porque dele advêm não apenas a
modernização da agricultura, mas a transferência de um modelo de
desenvolvimento econômico e de relações sociais para o Terceiro Mundo ─
o modelo \textbf{capitalista} de produção. Dessa forma, o
\textit{agribusiness} apenas \textbf{exacerba as desigualdades sociais}
que (\ldots) são as causas reais da fome.'' (Burbach \& Flynn, 1980,
\textit{apud} Pompeia, p.~79; grifo meu)

``De acordo com a análise {[}de Susan George em
\textit{O mercado da fome: as verdadeiras razões da fome no mundo}{]}, o
\textit{agribusiness} norte-americano
\textbf{destruía sistemas alimentares} e preferências dos consumidores ─
com o alastramento de produtos comestíveis e bebidas açucaradas que
\textbf{substituíam dietas tradicionais} mais saudáveis e ancoradas em
valores \textbf{locais} ─, além de
\textbf{desestabilizar padrões de emprego e estruturas comunitárias} nos
países em desenvolvimento.'' (Pompeia, p.~78; grifo meu)
\end{quote}

\hypertarget{contraproposta-uxe0-pl-do-veneno-pnara}{%
\subsection{Contraproposta à PL do Veneno:
PNARA}\label{contraproposta-uxe0-pl-do-veneno-pnara}}

!Pasted image 20220528170925.png !Pasted image 20220528170944.png
!Pasted image 20220528171015.png (Fonte:
\href{https://editora.redeunida.org.br/wp-content/uploads/2021/07/Livro-Dossie-–-Contra-o-Pacote-do-Veneno-e-em-Defesa-da-Vida.pdf}{``Dossiê:
Contra o Pacote do Veneno e em Defesa da Vida!''}, p.~40-42)

\begin{center}\rule{0.5\linewidth}{0.5pt}\end{center}

\hypertarget{referuxeancias}{%
\subsubsection{Referências}\label{referuxeancias}}

\hypertarget{livros}{%
\paragraph{Livros}\label{livros}}

\begin{itemize}
\tightlist
\item
  Karl Marx, \textbf{O Capital, Livro I}, Ed. Boitempo.
\item
  Karl Marx e Friedrich Engels, \textbf{A Ideologia Alemã}, Ed.
  Boitempo.
\item
  Ana Primavesi, \textbf{Cartilha da Terra}, Ed. Expressão Popular.
\item
  Caio Pompeia, \textbf{Formação Política do Agronegócio}, Ed. Elefante.
\end{itemize}

\hypertarget{publicauxe7uxf5es}{%
\paragraph{Publicações}\label{publicauxe7uxf5es}}

\begin{itemize}
\tightlist
\item
  \href{https://ojoioeotrigo.com.br/2021/08/agro-alimenta-o-mundo/}{O
  Agro brasileiro alimenta o mundo? Estudo da Embrapa usa regra de três
  para provar que sim, mas os fatos dizem que não (O Joio e o Trigo)}
\item
  Friedrich, Karen, et
  al.~\href{https://editora.redeunida.org.br/wp-content/uploads/2021/07/Livro-Dossie-–-Contra-o-Pacote-do-Veneno-e-em-Defesa-da-Vida.pdf}{``Dossiê:
  Contra o Pacote do Veneno e em Defesa da Vida!''}, 2021.

  \begin{itemize}
  \tightlist
  \item
    \href{https://actually-shoe-d2e.notion.site/Dossi-Contra-o-pacote-do-veneno-e-em-defesa-da-vida-bddf885061704f3b9a6fca367c1efb88}{Resumo
    feito pelo Fred Santiago (Ecoar-MG)} sobre o dossiê do Pacote do
    Veneno
  \end{itemize}
\item
  \textbf{Declaración final del foro mundial sobre soberanía alimentaria}
  (La Habana, Cuba, 7 de setembro de 2001). Disponível em:
  \url{https://base.socioeco.org/docs/doc-792_es.pdf}
\end{itemize}

\hypertarget{podcasts}{%
\paragraph{Podcasts}\label{podcasts}}

Série investigativa ``\textbf{Muito Além da Porteira}'', do podcast
\textit{Prato Cheio} (do blog \textit{O Joio e o Trigo}): - Ep. 1:
\href{https://ojoioeotrigo.com.br/2021/11/quando-a-faria-lima-encontra-a-boiada/}{Quando
a Faria Lima encontra a boiada} -
\href{https://ojoioeotrigo.com.br/2021/11/agronegocio-boom-inedito-mercado-financeiro/}{Exclusivo:
agronegócio vive boom inédito no mercado financeiro (O Joio e o Trigo)}
-
\href{https://ojoioeotrigo.com.br/2021/11/mercado-financeiro-e-agronegocio/}{O
passo a passo da união entre mercado financeiro e agronegócio (O Joio e
o Trigo)} - Ep 2:
\href{https://ojoioeotrigo.com.br/2021/11/venha-investir-na-sua-propria-fome/}{Venha
investir na sua própria fome} -
\href{https://ojoioeotrigo.com.br/2021/11/conheca-o-cassino-da-fome-em-wall-street-onde-os-donos-do-dinheiro-apostam-seu-prato/}{Conheça
o cassino da fome em Wall Street, onde os donos do dinheiro apostam seu
prato (O Joio e o Trigo)} -
\href{https://ojoioeotrigo.com.br/2021/11/blackrock-o-gigante-devorando-a-colheita/}{BlackRock,
o gigante devorando a colheita (O Joio e o Trigo)} - Ep 3:
\href{https://ojoioeotrigo.com.br/2022/01/ideias-para-acelerar-o-fim-do-mundo/}{Ideias
para acelerar o fim do mundo} -
\href{https://ojoioeotrigo.com.br/2021/11/planos-de-aposentadoria-nos-estados-unidos-alimentam-a-destruicao-da-floresta-amazonica/}{Planos
de aposentadoria nos Estados Unidos alimentam a destruição da Floresta
Amazônica (O Joio e o Trigo)} - Ep 4:
\href{https://ojoioeotrigo.com.br/2022/01/a-guerra-cultural-do-agronegocio/}{A
guerra cultural do agronegócio}\\
- Ep 5:
\href{https://ojoioeotrigo.com.br/2022/02/a-guerra-cultural-do-agronegocio-2/}{Calibã,
o agro e as bruxas}

\hypertarget{cursos}{%
\paragraph{Cursos}\label{cursos}}

\begin{itemize}
\tightlist
\item
  Curso
  \href{https://plataforma.laurocampos.org.br/cursos/crise-ambiental/}{``A
  crise ambiental e os desafios da esquerda brasileira''}, Fundação
  Lauro Campos e Marielle Franco (FLCMF-PSOL)

  \begin{itemize}
  \tightlist
  \item
    Módulo 3: \textbf{Amazônia e os desafios da Humanidade} (Adolfo
    Oliveira, UFPA)
  \item
    Módulo 6: \textbf{Soberania e Segurança Alimentar e Nutricional}
    (Anelise Rizzolo)
  \item
    Módulo 7: \textbf{Democratização da Terra} (Kelli Mafort, MST)
  \item
    Módulo 8: \textbf{Mineração e seus impatos sociais e ambientais}
    (Marcos Ferreira)
  \end{itemize}
\end{itemize}

\hypertarget{notuxedcias-sobre-desabastecimento-no-brasil}{%
\paragraph{Notícias sobre (Des)abastecimento no
Brasil}\label{notuxedcias-sobre-desabastecimento-no-brasil}}

\begin{itemize}
\tightlist
\item
  \href{https://ojoioeotrigo.com.br/2020/09/o-governo-deveria-estocar-arroz-nao-voce/}{O
  governo deveria estocar arroz, não você (O Joio e o Trigo, setembro
  2020)}
\item
  \href{https://www.bbc.com/portuguese/brasil-59070059}{Por que
  agricultores brasileiros estão deixando de plantar feijão --- e o que
  isso tem a ver com a fome (BBC News, novembro 2021)}
\end{itemize}

\hypertarget{vuxeddeos}{%
\paragraph{Vídeos}\label{vuxeddeos}}

\begin{itemize}
\tightlist
\item
  \href{https://www.youtube.com/watch?v=hHI61GHNGJM}{\textit{What is Solarpunk?}
  (Andrewism)}
\item
  \href{https://www.youtube.com/watch?v=lFDm7JLze5g}{\textit{A Quick Guide to Permablitzing}
  (Andrewism)}
\item
  \href{https://www.youtube.com/watch?v=u-JvyfZVkIM}{\textit{How we can Make Solarpunk a Reality}
  (Andrewism \& Our Changing Climate)}
\end{itemize}
